% GENERATED FILE. Edit canonical lessons, then run npm run generate:books.
% canonical-source-hash: fnv1a64:9f480798c8722e4b
% canonical-lessons: UR-C506-safai, UR-C506-manjhna, UR-C506-chipkana, UR-C506-dhulna, UR-C506-khangalna

\chapter{Scrubbing and Rinsing}
\label{ch:ur-scrubbing-and-rinsing}
\hlchaptermodality{506}

\begin{chapteropening}
\textbf{By the end of this chapter:} \emph{I can talk about cleaning the house and the dishes.}

Complete the last lesson of chapter 506: I can talk about cleaning the house and the dishes.
\end{chapteropening}

\section[\texorpdfstring{\ur{صفائی} (safāī)}{صفائی (safāī)}]{\ur{صفائی} (safāī) — cleaning; cleanliness}

\label{lesson:UR-C506-safai}



\begin{quote}
Before the new one: say the Urdu for to tear (something), then the Urdu for to have (something) put away.
\end{quote}

\subsection*{You'll want to know: \ur{صفائی}}
\textbf{\ur{صفائی}} — \emph{safāī} — \textquotedblleft{}cleaning; cleanliness\textquotedblright{}.

\begin{quote}
\ur{ہم} \ur{ہر} \ur{اتوار} \ur{کو} \ur{گھر} \ur{کی} \ur{صفائی} \ur{کرتے} \ur{ہیں۔}
\end{quote}

\begin{quote}
\emph{ham har itvār ko ghar kī safāī karte haiṅ.}
\end{quote}

\begin{quote}
We clean the house every Sunday.
\end{quote}

Feminine.

\begin{cousinweb}
\textbf{\ur{صفائی}} (\emph{safāī}) is built on \emph{sāf}, \textquotedblleft{}clean\textquotedblright{}, which is usually traced, through Persian, to Arabic \emph{ṣāfī}, \textquotedblleft{}pure, clear\textquotedblright{}, with the noun ending \emph{-āī}.
\end{cousinweb}

\begin{grammarlens}[title={What you've built}]
One more word for describing what you lived and felt.
\end{grammarlens}

\subsection*{Your turn}
\begin{itemize}
\raggedright
  \item \emph{Say it:} \emph{safāī}
  \item \emph{Say it:} \emph{safāī}, once more
  \item \emph{Say it:} \emph{ham har itvār ko ghar kī safāī karte haiṅ.} — We clean the house every Sunday.
  \item \emph{Recall:} say the Urdu for to tear (something), then the Urdu for to have (something) put away, then say \emph{safāī} again
\end{itemize}

\begin{tcolorbox}[breakable,colback=teal!4,colframe=teal!35!black,title={Before you move on}]
What does \ur{صفائی} mean? (\textquotedblleft{}Cleaning; cleanliness\textquotedblright{}.) Say it once more.
\end{tcolorbox}
\section[\texorpdfstring{\ur{مانجھنا} (māṅjhnā)}{مانجھنا (māṅjhnā)}]{\ur{مانجھنا} (māṅjhnā) — to scour (pots)}

\label{lesson:UR-C506-manjhna}



\begin{quote}
Before the new one: say the Urdu for to have (something) put away, then the Urdu for cleaning.
\end{quote}

\subsection*{You'll want to know: \ur{مانجھنا}}
\textbf{\ur{مانجھنا}} — \emph{māṅjhnā} — \textquotedblleft{}to scour (pots)\textquotedblright{}.

\begin{quote}
\ur{اس} \ur{نے} \ur{راکھ} \ur{سے} \ur{برتن} \ur{مانجھے۔}
\end{quote}

\begin{quote}
\emph{us ne rākh se bartan māṅjhe.}
\end{quote}

\begin{quote}
She scoured the pots with ash.
\end{quote}

Past, with \ur{نے} (he or she …): \ur{اس} \ur{نے} \ur{کڑاہی} \ur{مانجھی}, \emph{us ne kaṛāhī māṅjhī}.

\begin{cousinweb}
\textbf{\ur{مانجھنا}} (\emph{māṅjhnā}) is inherited: CDIAL lists the Hindi cousin \emph{mā̃jnā}, \textquotedblleft{}to rub, to clean, to polish\textquotedblright{}, under Sanskrit \emph{mārjati}, \textquotedblleft{}wipes, cleans\textquotedblright{} (CDIAL 10080).
\end{cousinweb}

\begin{grammarlens}[title={What you've built}]
One more word for describing what you lived and felt.
\end{grammarlens}

\subsection*{Your turn}
\begin{itemize}
\raggedright
  \item \emph{Say it:} \emph{māṅjhnā}
  \item \emph{Say it:} \emph{māṅjhnā}, once more
  \item \emph{Say it:} \emph{us ne rākh se bartan māṅjhe.} — She scoured the pots with ash.
  \item \emph{Recall:} say the Urdu for to have (something) put away, then the Urdu for cleaning, then say \emph{māṅjhnā} again
\end{itemize}

\begin{tcolorbox}[breakable,colback=teal!4,colframe=teal!35!black,title={Before you move on}]
What does \ur{مانجھنا} mean? (\textquotedblleft{}To scour (pots)\textquotedblright{}.) Say it once more.
\end{tcolorbox}
\section[\texorpdfstring{\ur{چپکانا} (chipkānā)}{چپکانا (chipkānā)}]{\ur{چپکانا} (chipkānā) — to stick (something) on; to glue}

\label{lesson:UR-C506-chipkana}



\begin{quote}
Before the new one: say the Urdu for cleaning, then the Urdu for to scour (pots).
\end{quote}

\subsection*{You'll want to know: \ur{چپکانا}}
\textbf{\ur{چپکانا}} — \emph{chipkānā} — \textquotedblleft{}to stick (something) on; to glue\textquotedblright{}.

\begin{quote}
\ur{میں} \ur{نے} \ur{لفافے} \ur{پر} \ur{ٹکٹ} \ur{چپکایا۔}
\end{quote}

\begin{quote}
\emph{maiṅ ne lifāfe par ṭikaṭ chipkāyā.}
\end{quote}

\begin{quote}
I stuck a stamp on the envelope.
\end{quote}

Past, with \ur{نے} (he or she …): \ur{اس} \ur{نے} \ur{دیوار} \ur{پر} \ur{اشتہار} \ur{چپکایا}, \emph{us ne dīvār par ishtihār chipkāyā}.

\begin{cousinweb}
\textbf{\ur{چپکانا}} (\emph{chipkānā}) is inherited: CDIAL lists the Hindi cousins \emph{cep}, \textquotedblleft{}gummy matter\textquotedblright{}, and \emph{cepnā}, \textquotedblleft{}to stick on\textquotedblright{}, under \emph{*cippa}, \textquotedblleft{}gummy matter\textquotedblright{} (CDIAL 4819).
\end{cousinweb}

\begin{grammarlens}[title={What you've built}]
One more word for describing what you lived and felt.
\end{grammarlens}

\subsection*{Your turn}
\begin{itemize}
\raggedright
  \item \emph{Say it:} \emph{chipkānā}
  \item \emph{Say it:} \emph{chipkānā}, once more
  \item \emph{Say it:} \emph{maiṅ ne lifāfe par ṭikaṭ chipkāyā.} — I stuck a stamp on the envelope.
  \item \emph{Recall:} say the Urdu for cleaning, then the Urdu for to scour (pots), then say \emph{chipkānā} again
\end{itemize}

\begin{tcolorbox}[breakable,colback=teal!4,colframe=teal!35!black,title={Before you move on}]
What does \ur{چپکانا} mean? (\textquotedblleft{}To stick (something) on; to glue\textquotedblright{}.) Say it once more.
\end{tcolorbox}
\section[\texorpdfstring{\ur{دھلنا} (dhulnā)}{دھلنا (dhulnā)}]{\ur{دھلنا} (dhulnā) — to be washed; to come clean}

\label{lesson:UR-C506-dhulna}



\begin{quote}
Before the new one: say the Urdu for to scour (pots), then the Urdu for to stick (something) on.
\end{quote}

\subsection*{You'll want to know: \ur{دھلنا}}
\textbf{\ur{دھلنا}} — \emph{dhulnā} — \textquotedblleft{}to be washed; to come clean\textquotedblright{}.

\begin{quote}
\ur{یہ} \ur{چادر} \ur{جلدی} \ur{دھل} \ur{گئی۔}
\end{quote}

\begin{quote}
\emph{yih chādar jaldī dhul gaī.}
\end{quote}

\begin{quote}
This sheet washed clean quickly.
\end{quote}

Past (it …): \ur{وہ} \ur{چادر} \ur{دھل} \ur{گئی}, \emph{voh chādar dhul gaī}.

\begin{cousinweb}
\textbf{\ur{دھلنا}} (\emph{dhulnā}) is the intransitive partner of \emph{dhonā}, \textquotedblleft{}to wash\textquotedblright{}, which CDIAL files under \emph{*dhauvati}, \textquotedblleft{}washes\textquotedblright{}, a form it reconstructs (CDIAL 6886).
\end{cousinweb}

\begin{grammarlens}[title={What you've built}]
One more word for describing what you lived and felt.
\end{grammarlens}

\subsection*{Your turn}
\begin{itemize}
\raggedright
  \item \emph{Say it:} \emph{dhulnā}
  \item \emph{Say it:} \emph{dhulnā}, once more
  \item \emph{Say it:} \emph{yih chādar jaldī dhul gaī.} — This sheet washed clean quickly.
  \item \emph{Recall:} say the Urdu for to scour (pots), then the Urdu for to stick (something) on, then say \emph{dhulnā} again
\end{itemize}

\begin{tcolorbox}[breakable,colback=teal!4,colframe=teal!35!black,title={Before you move on}]
What does \ur{دھلنا} mean? (\textquotedblleft{}To be washed; to come clean\textquotedblright{}.) Say it once more.
\end{tcolorbox}
\section[\texorpdfstring{\ur{کھنگالنا} (khaṅgālnā)}{کھنگالنا (khaṅgālnā)}]{\ur{کھنگالنا} (khaṅgālnā) — to rinse out}

\label{lesson:UR-C506-khangalna}



\begin{quote}
Before the new one: say the Urdu for to stick (something) on, then the Urdu for to be washed.
\end{quote}

\subsection*{You'll want to know: \ur{کھنگالنا}}
\textbf{\ur{کھنگالنا}} — \emph{khaṅgālnā} — \textquotedblleft{}to rinse out\textquotedblright{}.

\begin{quote}
\ur{اس} \ur{نے} \ur{نل} \ur{کے} \ur{نیچے} \ur{گلاس} \ur{کھنگالے۔}
\end{quote}

\begin{quote}
\emph{us ne nal ke nīche glās khaṅgāle.}
\end{quote}

\begin{quote}
She rinsed the glasses under the tap.
\end{quote}

Past, with \ur{نے} (he or she …): \ur{اس} \ur{نے} \ur{بالٹی} \ur{کھنگالی}, \emph{us ne balṭī khaṅgālī}.

\begin{cousinweb}
\textbf{\ur{کھنگالنا}} (\emph{khaṅgālnā}) is inherited: CDIAL lists the Hindi cousin \emph{khãgālnā}, \textquotedblleft{}to rinse\textquotedblright{}, under \emph{*khakkhālayati}, a form it reconstructs (CDIAL 3762).
\end{cousinweb}

\begin{grammarlens}[title={What you've built}]
One more word for describing what you lived and felt.
\end{grammarlens}

\subsection*{Your turn}
\begin{itemize}
\raggedright
  \item \emph{Say it:} \emph{khaṅgālnā}
  \item \emph{Say it:} \emph{khaṅgālnā}, once more
  \item \emph{Say it:} \emph{us ne nal ke nīche glās khaṅgāle.} — She rinsed the glasses under the tap.
  \item \emph{Recall:} say the Urdu for to stick (something) on, then the Urdu for to be washed, then say \emph{khaṅgālnā} again
\end{itemize}

\begin{tcolorbox}[breakable,colback=teal!4,colframe=teal!35!black,title={Before you move on}]
What does \ur{کھنگالنا} mean? (\textquotedblleft{}To rinse out\textquotedblright{}.) Say it once more.
\end{tcolorbox}
